(1) 将 $\eta=\eta_0e^{i(kx-\omega t)}$ 代入对流波动方程，得 $(-i\omega+iUk)^2=c^2(ik)^2$，即 $(\omega-Uk)^2=c^2k^2$。对非零 $k$ 取平方根，有 $\omega=(U+c)k$ 或 $\omega=(U-c)k$，符号表示相对流体的两种方向。
(2) 每支色散为严格线性关系，故实验室相速度和群速度都分别是 $v_\pm=U\pm c$。在随流体以速度 $U$ 移动的坐标系中，$\omega'=\omega-Uk$，因此两支群速为 $\pm c$；均匀对流只带来 Galilean 速度叠加。
(3) 固定正频率时 $k_+=\omega/(U+c)$、$k_-=\omega/(U-c)$，后一支在 $U<c$ 时具有负波数和负群速度，可逆流传播。$U=c$ 时第二支对给定非零正频率没有有限 $k$ 解，其实验室群速度为零；$U>c$ 时两支实验室群速度均为正，均不能向 $-x$ 传信。
(4) $U=0$ 时两支为 $\omega=\pm ck$，波以相同速率向左右传播。临界流速 $U=c$ 的“声学地平线”在本题是上游特征线停止后退的运动学界面；真正空间局域的地平线还需非均匀 $U(x)$、边界和流体背景稳定性分析，不能由均匀模型直接宣布黑洞辐射。
